% Abhakseite „Das kann ich“ – Zone nur im Gesamt (\mitzone)
%% TODO Ich-kann-Sätze: Platzhalter wie in den Aufgabendateien
\begin{abhakseite}
\ifdefined\mitzone
\abhakgruppe{Kennst du schon}
\abhak{1}{Quadrieren auch negativer Zahlen und Dezimalzahlen, Quadrat vor Punkt vor Strich}
\abhak{2}{Quadratwurzel ziehen}
\abhak{3}{Lineare Gleichung mit Äquivalenzumformungen und Strich lösen, auch negative Vorzahl und x auf beiden Seiten}
\abhak{4}{Lösung durch Einsetzen prüfen mit (wA)/(fA), auch mit negativer Zahl}
\abhak{5}{Klammer zuerst und Punkt vor Strich mit negativen Zahlen}
\abhak{6}{Scheitelpunktform lesen}
\abhak{7}{Ausklammern, Ausmultiplizieren und binomische Formeln}
\abhak{8}{Terme ordnen und zusammenfassen, x² zuerst, dann x-Glieder, dann Zahlen}
\abhak{9}{Quadrieren auch negativer Zahlen und Dezimalzahlen, Quadrat vor Punkt vor Strich – Fehler finden}
\abhak{10}{Quadrieren auch negativer Zahlen und Dezimalzahlen, Quadrat vor Punkt vor Strich – selbst rechnen}
\fi
\abhakgruppe{Einheit 1 · Wurzelziehen und Lösbarkeit}
\abhak{11}{Quadratisch oder linear?}
\abhak{12}{Wurzelziehen und Lösbarkeit}
\abhak{13}{Wurzelziehen und Lösbarkeit – Fehler finden, Begründen, Darstellungswechsel, Anwendung}
\abhakgruppe{Einheit 2 · Satz vom Nullprodukt}
\abhak{14}{Nullprodukt}
\abhak{15}{Nullprodukt – Fehler finden, Begründen}
\abhakgruppe{Einheit 3 · Normalform und p-q-Formel}
\abhak{16}{Welche Form?}
\abhak{17}{Ist die Vorzahl von x² eins?}
\abhak{18}{p-q-Formel}
\abhak{19}{Gerade und Parabel gleichsetzen als Anwendung}
\abhak{20}{p-q-Formel – Fehler finden, Begründen, Darstellungswechsel, Anwendung}
\abhakgruppe{Einheit 4 · Sachaufgaben}
\abhak{21}{Sachaufgaben}
\abhak{22}{Quadrat mit Rand}
\abhak{23}{Sachaufgaben – Fehler finden, Begründen}
\end{abhakseite}
